\documentclass{article}
\usepackage{cancel}
\usepackage{gensymb}
\usepackage{tikz}
\usepackage{amsmath}
\usepackage{ amssymb }
\usepackage{listings}
\usepackage{geometry}
\usepackage{graphicx}
\usepackage{amsmath}
\usepackage{amsthm}
\usepackage{empheq}
\usepackage{mdframed}
\usepackage{booktabs}
\usepackage{lipsum}
\usepackage{color}
\usepackage{psfrag}
\usepackage{pgfplots}
\usepackage{bm}
\usepackage{xcolor}
\usepackage{minted}       
\usepackage{verbatim}    
\usepackage{siunitx}
\usepackage{hyperref}
\usepackage[
  backend=biber,
  sorting=none,
%   style=mla
]{biblatex}

\addbibresource{sources.bib} 

\usepgfplotslibrary{external}
\tikzexternalize


\usetikzlibrary{angles,quotes}
\usetikzlibrary{arrows.meta,calc}

% --- visual style for the cell ---
\newmdenv[
  backgroundcolor=gray!5,
  linecolor=gray!50,
  roundcorner=5pt,
  innerleftmargin=10pt,
  innerrightmargin=10pt,
  innertopmargin=10pt,
  innerbottommargin=10pt
]{codecell}

\newif\ifcodefast
\codefasttrue % change whether to write code fancy or just plain

\ifcodefast
  % ---------- FAST MODE ----------
  \newenvironment{codeblockPython}
    {\VerbatimEnvironment
     \begin{codecell}
     \begin{Verbatim}}
    {\end{Verbatim}
     \end{codecell}}
\else
  % ---------- FULL MINTED MODE ----------
  \newenvironment{codeblockPython}
    {\VerbatimEnvironment
     \begin{codecell}
     \begin{minted}{python3}}
    {\end{minted}
     \end{codecell}}
\fi

\geometry{a4paper}
\renewcommand{\thesubsubsection}{\alph{subsubsection})}
\setlength{\parindent}{0pt}

\DeclareSIUnit{\Msun}{M_{\odot}}
\DeclareSIUnit{\Lsun}{L_{\odot}}
\DeclareSIUnit{\Mpc}{Mpc}


\def\hwNUM{12}


\title{Astro HW \hwNUM}
\author{Pierson Lipschultz}

\begin{document}
\maketitle
\setcounter{section}{\hwNUM}

\subsection{Eddington luminosity and massive stars}

\subsubsection{}
Given \dots
\[L_{\rm Edd} \simeq 3.3\times10^4 \left(\frac{M}{M_\odot}\right) L_\odot\]

\[L \simeq L_\odot \left(\frac{M}{M_\odot}\right)^3\] 

\[L = L_{\rm Edd}, \qquad \text{when disruption occurs}\] 

We can derive \dots 
\[\Rightarrow\
L_\odot \left(\frac{M}{M_\odot}\right)^3
=
3.3\times10^4 \left(\frac{M}{M_\odot}\right) L_\odot\]

\[\Rightarrow
\left(\frac{M}{M_\odot}\right)^2 = 3.3\times10^4\]

\[\Rightarrow
M_{\max} \approx \sqrt{3.3\times10^4}\,M_\odot \approx 1.8\times10^2\,M_\odot
\]

\subsubsection{}
This calculated value of \(M_{lim} \approx 1.8 \times 10^2 M_\odot\) is much lower than a considerable amount of observed systems. For example, the current most massive star has a proposed mass of \(346\pm42 M_\odot\)\cite{Keszthelyi_2025}.

\subsection{Accretion disk temperatures}
\subsubsection{}
Given \dots 
\[T \approx \left[\frac{GM}{2\sigma_{SB}r^2 }\Sigma v\right]^{1/4}\]
\[\dot{M} = 2\pi r\Sigma(r)v(r)\]

We find \dots 
\[\dot{M} = 2\pi r\Sigma(r)v(r)\]
\[\Sigma(r)v = \frac{\dot{M}}{r2\pi}\]
\[T \approx \left[\frac{GM}{2\sigma_{SB}r^2 }\frac{\dot{M}}{r2\pi}\right]^{1/4}\]
\[T \approx \left[\frac{GM\dot{M}}{4\pi\sigma_{SB}r^2 }\right]^{1/4}\]

\subsubsection{}
Given \dots
\[T \approx \left[\frac{GM\dot{M}}{4\pi\sigma_{SB}r^2 }\right]^{1/4}\]
\[L = \eta \dot{M}c^2\]
\[L = E_L, \quad \eta \sim .1, \quad r = \frac{6GM}{c^2}\]

\subsubsection{}

\printbibliography[
heading=bibintoc,
title={\centering Sources}
]

\end{document}